\chapter{Cómo usar este manual}

\opening{Un helicóptero liviano se vuela con la cabeza antes que con las manos: lo que se estudia en tierra es lo que responde por uno en el aire.}

El Robinson R22 es un helicóptero biplaza, monomotor, con rotor principal bipala de tipo balancín (\emph{teetering}) y tren de patines. Es el entrenador más difundido del mundo y, al mismo tiempo, una máquina con un margen de error estrecho: rotor de baja inercia, controles muy sensibles y limitaciones que deben conocerse de memoria. Este libro reúne en un solo lugar, en español y ordenado por temas, toda la información que el fabricante y la guía de prevuelo entregan al piloto y al mecánico.

\section{Organización}

El libro está dividido en cuatro partes y dos apéndices:

\begin{itemize}
  \item \textbf{Parte I --- Generalidades.} Datos descriptivos del helicóptero, dimensiones, motorizaciones, definiciones de performance y de peso y balance.
  \item \textbf{Parte II --- Operación.} El contenido completo del Manual de Operación del Piloto (POH): limitaciones, procedimientos de emergencia y normales, performance, peso y balance, descripción de sistemas, manejo y servicio, suplementos (versión policial, Mariner y aviónica opcional) y la sección de seguridad con el texto íntegro de todas las \emph{Safety Notices} vigentes.
  \item \textbf{Parte III --- Inspección prevuelo.} La lista de chequeo diaria del fabricante, desarrollada ítem por ítem con la explicación de qué mirar y por qué, junto con los casos de investigación de accidentes que ilustran cada punto, y una versión compacta para usar junto al helicóptero.
  \item \textbf{Parte IV --- Mantenimiento.} Una síntesis organizada del Manual de Mantenimiento e Instrucciones de Aeronavegabilidad Continuada (MM): programa de inspecciones, componentes de vida limitada, especificaciones de servicio, límites de desgaste, valores de torque, reglaje y procedimientos clave de cada sistema.
  \item \textbf{Apéndices.} Abreviaturas y glosario bilingüe, y tablas de conversión.
\end{itemize}

\section{Convenciones}

El fabricante usa tres tipos de aviso, que se conservan con la siguiente presentación:

\begin{advertencia}
Indica que puede producirse una lesión personal o la muerte si no se sigue la instrucción. (En el MM, \emph{WARNING}.)
\end{advertencia}

\begin{precaucion}
Indica que pueden producirse daños al equipo, lesiones o la muerte si no se sigue el procedimiento o la instrucción. (\emph{CAUTION} en el original.)
\end{precaucion}

\begin{nota}
Entrega énfasis o información complementaria. (\emph{NOTE} en el original.)
\end{nota}

Los recuadros \textbf{Caso de investigación} de la Parte III resumen accidentes e incidentes reales, investigados por la NTSB, la ATSB y otras autoridades, que ilustran la consecuencia de omitir el ítem que se está revisando. Muchos involucran otros modelos de helicóptero: lo que importa es la lección, no la marca.

Las listas de chequeo se presentan con el ítem a la izquierda y la condición o acción esperada a la derecha, igual que en el original:

\begin{checklist}
\chk{Interruptor de batería}{ON}
\chk{Holgura de correas en V}{1,5 in (4 cm) máx.}
\end{checklist}

\section{Terminología}

Se ha preferido la terminología aeronáutica usada en Chile y Latinoamérica. Cuando un término inglés es de uso habitual en la cabina o en el taller (\emph{Telatemp}, \emph{gascolator}, \emph{sprag}, \emph{low-G}) se conserva entre paréntesis o en cursiva la primera vez. Los nombres de luces e interruptores se mantienen tal como están rotulados en el panel (por ejemplo, la luz \textsf{LOW RPM} o el interruptor \textsf{CLUTCH}), porque es así como el piloto los verá. El apéndice~\ref{ap:abreviaturas} contiene un glosario inglés--español.

\section{Variantes del R22}

El certificado tipo FAA H10WE cubre las versiones R22 (Standard), R22 HP, R22 Alpha, R22 Beta, R22 Beta II y las variantes Mariner (con flotadores) y policial. Muchas limitaciones dependen de la variante y del motor instalado; cuando corresponde, las tablas separan los valores por modelo. Antes de usar cualquier dato, identifique con certeza la variante, el motor y el equipamiento opcional de su helicóptero.
